% Môn: Vật lý
% Lớp: 12
% Chương: Chương I. Vật lí nhiệt
% Bài: L12C1 Bài 6. Nhiệt hóa hơi riêng · __CHUA_PHAN_LOAI__
% ===== Câu 3 =====
% ID: L12C1B6-196-TLN
% Mức: TH
\dangbt{__CHUA_PHAN_LOAI__}
\begin{ex}
% ID: L12C1B6-196-TLN
% Mức: TH
Cho hằng số Boltzmann $k=1,38.10^{-23}\ J/K$. Động năng tịnh tiến trung bình của phân tử khí lí tưởng ở $30^\circ C$ có giá trị là
\shortans{A}
\loigiai{
• $\overline{E_\text{đ}}=\dfrac{3}{2}kT=\dfrac{3}{2}.1,38.10^{-23}.(30+273)=6,272.10^{-21}\ (J)$
}
\end{ex}

% ===== Câu 7 =====
% ID: L12C1B6-197-TLN
% Mức: TH
\begin{ex}
% ID: L12C1B6-197-TLN
% Mức: TH
Một bình chứa khí lí tưởng neon (Ne) có khối lượng riêng là $1,4\,\text{kg}/m^3$. Biết căn bậc hai giá trị trung bình của các bình phương tốc độ phân tử khí $\left(\sqrt{\overline{v^2}}\right)$ là $450\,\text{m/s}$. Áp suất khí trong bình là
\shortans{A}
\loigiai{
Áp suất khí trong bình $p=\dfrac{1}{3}\dfrac{Nm}{V}\overline{v^2}=\dfrac{1}{3}\rho\overline{v^2}=\dfrac{1}{3}.1,4.450^2=94\,500\ Pa$
}
\end{ex}

% ===== Câu 11 =====
% ID: L12C1B6-198-TLN
% Mức: VD
\begin{ex}
% ID: L12C1B6-198-TLN
% Mức: VD
Một bình chứa khối khí có khối lượng riêng $\rho$, $\overline{v^2}$ là trung bình của các bình phương tốc độ phân tử. Áp suất khí trong bình theo mô hình động học phân tử là
\shortans{A}
\loigiai{
Áp suất khí trong bình theo mô hình động học phân tử là $p=\dfrac{1}{3} \rho \overline{v^2}$.
}
\end{ex}

% ===== Câu 15 =====
% ID: L12C1B6-199-TLN
% Mức: VD
\begin{ex}
% ID: L12C1B6-199-TLN
% Mức: VD
Số phân tử khí hydrogen chứa trong $1\,\mathrm{m^3}$ có áp suất $200\,\mathrm{mmHg}$ và tốc độ căn quân phương $2400\,\mathrm{m/s}$ là
\shortans{A}
\loigiai{
Với hydrogen, khối lượng một phân tử là $m_0=\dfrac{2\cdot 10^{-3}}{6,02\cdot 10^{23}}\approx 3,32\cdot 10^{-27}\,\mathrm{kg}.$ Ta có $p=\dfrac{1}{3}\mu m_0v^2\Rightarrow \mu=\dfrac{3p}{m_0v^2}.$ Vì $V=1\,\mathrm{m^3}$ nên $N=\mu V=\dfrac{3\cdot 200\cdot 101325}{760\cdot 3,32\cdot 10^{-27}\cdot 2400^2}\approx4,2\cdot 10^{24}\ \text{phân tử}.$
}
\end{ex}

% ===== Câu 19 =====
% ID: L12C1B6-200-TLN
% Mức: VD
\begin{ex}
% ID: L12C1B6-200-TLN
% Mức: VD
Khối lượng riêng của một chất khí bằng $6\cdot 10^{-2}\,\mathrm{kg/m^3}$, tốc độ căn quân phương của chúng là $500\,\mathrm{m/s}$. Áp suất mà khí đó tác dụng lên thành bình là
\shortans{D}
\loigiai{
Ta có $p=\dfrac{1}{3}\rho v^2=\dfrac{1}{3}\cdot 6\cdot 10^{-2}\cdot 500^2=5\cdot 10^3\,\mathrm{Pa}.$
}
\end{ex}

% ===== Câu 23 =====
% ID: L12C1B6-201-TLN
% Mức: VD
\begin{ex}
% ID: L12C1B6-201-TLN
% Mức: VD
Khối lượng phân tử khí hydrogen là $3,3\cdot 10^{-24}\,\mathrm{g}$. Biết rằng trong $1\,\mathrm{s}$, có $10^{23}$ phân tử khí hydrogen chuyển động với vận tốc $1000\,\mathrm{m/s}$ đập vào $1\,\mathrm{cm^2}$ thành bình theo phương nghiêng $30^\circ$ với thành bình. Áp suất khí lên thành bình theo đơn vị $10^3\,\mathrm{N/m^2}$ là bao nhiêu?
\shortans{3{,}3}
\loigiai{
Khối lượng một phân tử hydrogen: $m=3,3\cdot 10^{-24}\,\mathrm{g}=3,3\cdot 10^{-27}\,\mathrm{kg}.$ Với một phân tử đập vào thành bình theo phương nghiêng $30^\circ$ với thành bình, biến thiên động lượng theo phương vuông góc thành bình là $\Delta p=2mv\sin30^\circ=mv.$ Trong $\Delta t=1\,\mathrm{s}$, có $N=10^{23}$ phân tử đập vào diện tích $S=1\,\mathrm{cm^2}=10^{-4}\,\mathrm{m^2}$, nên áp suất là $p=\dfrac{Nmv}{S\Delta t}=\dfrac{10^{23}\cdot 3,3\cdot 10^{-27}\cdot 10^3}{10^{-4}\cdot 1}=3,3\cdot 10^3\,\mathrm{N/m^2}.$ Vậy đáp số là $3,3$.
}
\end{ex}

% ===== Câu 27 =====
% ID: L12C1B6-202-TLN
% Mức: VD
\begin{ex}
% ID: L12C1B6-202-TLN
% Mức: VD
Xét khối khí chứa trong một bình kín, biết mật độ động năng phân tử, tức tổng động năng tịnh tiến trung bình của các phân tử khí trong $1\,\mathrm{m^3}$ thể tích khí, có giá trị $10^{-4}\,\mathrm{J/m^3}$. Tính áp suất của khí trong bình, lấy đơn vị $10^{-5}\,\mathrm{Pa}$ và làm tròn đến $2$ chữ số sau dấu phẩy thập phân.
\shortans{6{,}67}
\loigiai{
Gọi $w=10^{-4}\,\mathrm{J/m^3}$ là mật độ động năng. Ta có $p=\dfrac{2}{3}w=\dfrac{2\cdot 10^{-4}}{3}=6,67\cdot 10^{-5}\,\mathrm{Pa}.$ Vậy đáp số theo đơn vị $10^{-5}\,\mathrm{Pa}$ là $6,67$.
}
\end{ex}

% ===== Câu 31 =====
% ID: L12C1B6-203-TLN
% Mức: VD
\begin{ex}
% ID: L12C1B6-203-TLN
% Mức: VD
Một bình dung tích $7,5\,\mathrm{L}$ chứa $24\,\mathrm{g}$ khí oxygen ở áp suất $2,5\cdot 10^5\,\mathrm{N/m^2}$. Động năng trung bình của các phân tử khí oxygen theo đơn vị $10^{-21}\,\mathrm{J}$ là bao nhiêu?
\shortans{6,23}
\loigiai{
Ta có $p=\dfrac{2}{3}\mu\overline{W}_{\mathrm{đ}}\Rightarrow \overline{W}_{\mathrm{đ}}=\dfrac{3p}{2\mu}.$ Mặt khác $\mu=\dfrac{N}{V}=\dfrac{mN_A}{MV}.$ Do đó $\overline{W}_{\mathrm{đ}}=\dfrac{3pMV}{2mN_A}=\dfrac{3\cdot 2,5\cdot 10^5\cdot 32\cdot 7,5\cdot 10^{-3}}{2\cdot 24\cdot 6,02\cdot 10^{23}}\approx6,23\cdot 10^{-21}\,\mathrm{J}.$ Vậy đáp số là $6,23$.
}
\end{ex}

% ===== Câu 35 =====
% ID: L12C1B6-204-TLN
% Mức: VD
\begin{ex}
% ID: L12C1B6-204-TLN
% Mức: VD
Tính tốc độ trung bình của bình phương vận tốc của các phân tử không khí (làm tròn kết quả đến chữ số hàng đơn vị)?
\shortans{$494$}
\loigiai{
• Tốc độ trung bình của bình phương vận tốc của các phân tử không khí là: $\sqrt{\overline {v^2}} =\sqrt{\dfrac{3.p}{\rho}}=\sqrt{\dfrac{3.1,01.10^5}{1,24}}=494\ m/s$
}
\end{ex}

% ===== Câu 39 =====
% ID: L12C1B6-205-TLN
% Mức: VD
\begin{ex}
% ID: L12C1B6-205-TLN
% Mức: VD
Một hộp kín hình lập phương cạnh $10\,\mathrm{cm}$ chứa khí lí tưởng đơn nguyên tử ở nhiệt độ $20^\circ\mathrm{C}$ và áp suất $1,2\cdot 10^6\,\mathrm{Pa}$. Số phân tử khí chứa trong bình là
\shortans{A}
\loigiai{
Cạnh hình lập phương là
$a=10 \mathrm{cm}=0,1 \mathrm{m}.$
Thể tích bình:
$V=a^3=(0,1)^3=10^{-3} \mathrm{m^3}.$
Ta có
$T=20+273=293 \mathrm{K}.$
Áp dụng phương trình khí lí tưởng theo số phân tử:
$pV=NkT.$
Suy ra
$N=\dfrac{pV}{kT}.$
Thay số:
 N= $\dfrac{1,2\cdot 10^6\cdot 10^{-3}}$ {1,38 $\cdot$ 10^{-23} $\cdot$ 293}
$\approx$ 2,97 $\cdot$ 10^{23}.
}
\end{ex}
